% Part: first-order-logic
% Chapter: models-theories
% Section: theories

\documentclass[../../../include/open-logic-section]{subfiles}

\begin{document}

\olfileid[tr]{fol}{mat}{the}

\olsection{Birinci Dereceden Kuram Örnekleri}

\begin{ex}
$\Lang L_<$ dilindeki katı doğrusal sıralamalar kuramı
\begin{align*}
\{\quad & \lforall[x][\lnot x < x], \\
& \lforall[x][\lforall[y][((x < y \lor y <
    x) \lor x = y)]], \\
& \lforall[x][\lforall[y][\lforall[z][((x < y
      \land y < z) \lif x < z)]]] \quad \}
\end{align*}
kümesiyle aksiyomlaştırılır. Bu küme amaçlanan !!{structure}s
ifadelerini tam olarak yakalar: Her katı doğrusal sıralama bu aksiyom
sisteminin bir modelidir; tersine, $R$, $X$ kümesi üzerinde bir
doğrusal sıralamaysa $\Domain M=X$ ve $\Assign{<}{M}=R$ olan
$\Struct M$ yapısı bu kuramın bir modelidir.
\end{ex}

\begin{ex}
$\Obj 1$ (!!{constant}) ve $\cdot$ (iki yerli !!{function}) simgelerini
içeren dilde gruplar kuramı
\begin{align*}
& \lforall[x][\eq[(x \cdot \Obj 1)][x]]\\
& \lforall[x][\lforall[y][\lforall[z][\eq[(x \cdot (y \cdot z))][((x
          \cdot y) \cdot z)]]]]\\
& \lforall[x][\lexists[y][\eq[(x \cdot y)][\Obj 1]]]
\end{align*}
ile aksiyomlaştırılır.
\end{ex}

\begin{ex}
Peano aritmetiği kuramı, aritmetik dili~$\Lang L_A$ içindeki aşağıdaki
tümcelerle aksiyomlaştırılır:
\begin{align*}
& \lforall[x][\lforall[y][(\eq[x'][y'] \lif \eq[x][y])]]\\
& \lforall[x][\eq/[\Obj 0][x']]\\
& \lforall[x][\eq[(x + \Obj 0)][x]]\\
& \lforall[x][\lforall[y][\eq[(x + y')][(x + y)']]]\\
& \lforall[x][\eq[(x \times \Obj 0)][\Obj 0]]\\
& \lforall[x][\lforall[y][\eq[(x \times y')][((x \times y) + x)]] ]\\
& \lforall[x][\lforall[y][(x < y \liff \lexists[z][\eq[(z' + x)][y])]]]
\intertext{ve ayrıca aşağıdaki biçimdeki bütün tümceler:}
& (!A(\Obj 0) \land \lforall[x][(!A(x) \lif !A(x'))]) \lif \lforall[x][!A(x)].
\end{align*}
Son biçimde sonsuz sayıda tümce bulunduğu için bu aksiyom sistemi
sonsuzdur. Son biçime \emph{tümevarım şeması} denir. (Aslında tümevarım
şeması burada gösterdiğimizden biraz daha karmaşıktır.)

Son aksiyom,~$<$ bağıntısının \emph{açık bir tanımıdır}.
\end{ex}

\begin{ex}
Saf kümeler kuramı, matematiğin temellerinde (ve felsefesinde) önemli
bir rol oynar. Bir kümenin bütün !!{element}s yine saf kümelerse o
kümeye saf denir. Bu nedenle boş küme saftır; fakat !!a{element}
olarak küme olmayan bir şeyi içeren bir küme saf değildir. Dolayısıyla
saf kümeler, yalnızca boş kümeden ve kendileri küme olmayan nesneler
olan hiçbir ``temel elemandan'' yararlanmadan oluşturulan kümelerdir.

Aşağıdaki sistem, saf kümeler kuramı için bir aksiyom sistemi olarak
düşünülebilir:
\begin{align*}
& \lexists[x][\lnot \lexists[y][y \in x]]\\
& \lforall[x][\lforall[y][(\lforall[z](z \in x \liff z \in y) \lif
\eq[x][y])]]\\
& \lforall[x][\lforall[y][\lexists[z][\lforall[u][(u \in z \liff
(\eq[u][x] \lor \eq[u][y]))]]]]\\
& \lforall[x][\lexists[y][\lforall[z][(z \in y \liff \lexists[u][(z \in
        u \land u \in x)])]]]\\
\intertext{ve ayrıca aşağıdaki biçimdeki bütün tümceler:} &
\lexists[x][\lforall[y][(y \in x \liff !A(y))]].
\end{align*}
Birinci aksiyom hiç !!{element}s bulunmayan bir kümenin var olduğunu
(yani $\emptyset$ kümesinin varlığını); ikincisi kümelerin kaplamsal
olduğunu; üçüncüsü herhangi $X$ ve $Y$ kümeleri için $\{X,Y\}$
kümesinin var olduğunu; dördüncüsü ise herhangi bir $X$ kümesi için
$\cup X$ kümesinin var olduğunu söyler. Burada $\cup X$, $X$'in bütün
elemanlarının birleşimidir.

Son olarak anılan !!{sentence}s topluca \emph{naif küme oluşturma
şeması} diye adlandırılır. Bu şema özünde her $!A(x)$ için
$\Setabs{x}{!A(x)}$ kümesinin var olduğunu söyler; ilk bakışta doğru,
yararlı, hatta zorunlu görünebilecek bir aksiyomdur. ``Naif'' denmesinin
nedeni, bu kuramı sağlanamaz kıldığının ortaya çıkmasıdır: $!A(y)$
yerine $\lnot y\in y$ alırsak
\[
\lexists[x][\lforall[y][(y \in x \liff \lnot y \in y)]]
\]
!!{sentence} ifadesini elde ederiz ve bu !!{sentence} hiçbir
!!{structure} içinde sağlanmaz.
\end{ex}

\begin{ex}
\emph{Mereoloji} alanında \emph{parça olma} bağıntısı temel bir
bağıntıdır. Küme kuramlarında olduğu gibi, bu bağıntıya ilişkin çeşitli
(bazen birbiriyle çatışan) anlayışları aksiyomlaştıran parça olma
kuramları vardır.

Mereoloji dili tek bir iki yerli yüklem simgesi~$\Obj P$ içerir ve
$\Atom{\Obj P}{x,y}$, ``$x$, $y$'nin bir parçasıdır'' anlamına gelir.
Bu yorumu göz önünde tuttuğumuzda bu dil için !!a{structure}
ifadesine \emph{parça olma yapısı} denir. Elbette tek bir iki yerli
yüklem içeren her yapı bu adı gerçekten hak etmez. ``Parça olma''
anlamını yakalayabilmesi için $\Assign{\Obj P}{M}$ bazı koşulları
sağlamalıdır; bunları bir parça olma kuramının aksiyomları olarak
belirtebiliriz. Örneğin parça olma nesneler üzerinde bir kısmi
sıralamadır: Her nesne kendisinin (gerçi \emph{öz olmayan}) bir
parçasıdır; iki farklı nesne birbirinin parçası olamaz; bir nesnenin
bir parçasının parçası da o nesnenin parçasıdır. Bu anlamda ``parçası
olmak'', ``alt kümesi olmak'' bağıntısına benzer; fakat ne yansımalı ne
de geçişli olan ``elemanı olmak'' bağıntısına benzemez.
\begin{align*}
& \lforall[x][\Atom{\Obj P}{x,x}] \\
& \lforall[x][\lforall[y][((\Part{x}{y} \land \Part{y}{x})
      \lif \eq[x][y])]] \\
& \lforall[x][\lforall[y][\lforall[z][((\Part{x}{y} \land
        \Part{y}{z}) \lif \Part{x}{z})]]]\\
\intertext{Ayrıca herhangi iki nesnenin mereolojik bir toplamı vardır
(bu iki nesneyi parça olarak içeren ve bu bakımdan en küçük olan bir nesne):} &
\lforall[x][\lforall[y][\lexists[z][\lforall[u][(\Part{z}{u} \liff
        (\Part{x}{u} \land \Part{y}{u}))]]]].
\end{align*}
Bunlar, metafizikçilerin ele aldığı temel parça olma ilkelerinden
yalnızca bazılarıdır. Bununla birlikte, önce bazı bağıntılar
tanımlanmadan daha ileri ilkeleri formülleştirmek ya da yazmak hızla
zorlaşır. Örneğin mereolojiyle ilgilenen metafizikçilerin çoğu şu
ilkeyi de geçerli sayar: Bir $x$ nesnesinin bir öz parçası~$y$ varsa,
$x$ ayrıca $y$ ile hiçbir ortak parçası bulunmayan ve $y$ ile
birleşimi $x$ olan bir $z$ parçasına sahiptir.
\end{ex}

\end{document}
